% !Mode:: "TeX:UTF-8"

\chapter*{\centering \sanhao\bfseries ABSTRACT}

\xiaosi
\setstretch{1.5}

This report studies the cognitive boundaries of large language models in a controlled spatial reasoning task. Taking the wheelchair turning problem in a narrow stair corner as the fixed scenario, the project constructs prompts with different description levels and irrelevant perturbations, collects model responses, and extracts behavioral features such as final judgment, reasoning steps, formula usage, coordinate-system usage, and error category.

The analysis framework combines descriptive accuracy statistics, level consistency, noise-induced flip rate, decision tree analysis, K-Means clustering, and FP-Growth association rule mining. It aims to answer whether formal geometric descriptions improve model reliability, whether irrelevant environmental information changes model judgments, and what failure patterns appear in different models. This version provides a structured report framework based on the current project code and output artifacts. Final numerical findings can be added after multi-model response collection is completed.

\vspace{0.8cm}

\noindent {\bfseries\xiaosi Key words:} Large language model, Spatial reasoning, Data mining, Prompt engineering, Association rule mining
